% !TEX program = pdflatex
% !TEX encoding = UTF-8 Unicode
% =============================================================================
% Artículo para Tecnología en Marcha (TEC).
%
% Formato tomado de "Plantilla artículos Tecnología en marcha-OTH.docx":
% carta, márgenes de 2,5 cm arriba y abajo y 3 cm a los lados, texto a 12 pt,
% título en español y en inglés, autores con afiliación, correo y ORCID,
% resumen de 150 a 250 palabras, 5 a 7 palabras clave separadas por punto y
% coma, secciones sin numerar, "Cuadro N." arriba del cuadro, "Figura N."
% abajo de la figura, "Fuente:" en cada uno y referencias IEEE numeradas.
%
% Compilar en Overleaf con pdfLaTeX; Overleaf corre biber solo.
% =============================================================================
\documentclass[12pt,letterpaper]{article}

\usepackage[T1]{fontenc}
\usepackage[utf8]{inputenc}
\usepackage{newtxtext,newtxmath}
% Sin es-tabla: esa opción cambia "Cuadro" por "Tabla", y la plantilla usa
% "Cuadro".
\usepackage[spanish,es-noshorthands]{babel}
\usepackage[spanish=mexican]{csquotes}
\usepackage[letterpaper,top=2.5cm,bottom=2.5cm,left=3cm,right=3cm]{geometry}
\usepackage{graphicx}
\usepackage{booktabs}
\usepackage{tabularx}
\usepackage{array}
\usepackage{xcolor}
\usepackage{titlesec}
\usepackage{caption}
\usepackage[style=ieee,backend=biber]{biblatex}
\addbibresource{referencias.bib}
% biblatex-ieee pone la coma dentro de las comillas, como en inglés. La
% plantilla la deja afuera, como en español: “Título”, Revista.
\AtBeginBibliography{\stdpunctuation}
% Holgura para partir líneas en la bibliografía, donde abundan títulos largos
% sin puntos de corte (evita que una referencia se salga del margen).
\AtBeginBibliography{\setlength{\emergencystretch}{1.5em}}
% Abreviaturas de la plantilla ("vol. 5, no. 1, pp. 1-9") en lugar de las de
% biblatex en español ("n.º", "págs.").
\DefineBibliographyStrings{spanish}{%
  number = {no\adddot},
  pages  = {pp\adddot},
  page   = {p\adddot},
  url    = {Disponible en},
  % El estilo IEEE no traduce el número de artículo al español: sin esto,
  % biblatex imprime la clave interna ("artno") en negrita.
  artno  = {art\adddot\ no\adddot},
}
\usepackage[hidelinks]{hyperref}

% !TEX root = ../tecnologia-en-marcha.tex
% !TEX encoding = UTF-8 Unicode
% =============================================================================
% Marcas de revisión del borrador.
%
%   \pendiente{...}  algo que falta escribir o decidir (rojo)
%   \revisar{...}    una afirmación del texto que hay que verificar (naranja)
%
% Antes del envío, cambiar \notastrue por \notasfalse: las marcas desaparecen
% del PDF sin tocar el texto.
% =============================================================================
\newif\ifnotas
\notastrue

\newcommand{\pendiente}[1]{%
  \ifnotas{\color{red!70!black}\textbf{[Pendiente:} #1\textbf{]}}\fi}
\newcommand{\revisar}[1]{%
  \ifnotas{\color{orange!85!black}\textbf{[Revisar:} #1\textbf{]}}\fi}

% "Fuente:" debajo de cuadros, como pide la plantilla.
\newcommand{\fuente}[1]{\par\smallskip{\centering\small Fuente: #1\par}}


% --- Secciones sin numerar, como en la plantilla -----------------------------
\setcounter{secnumdepth}{0}
\titleformat{\section}{\normalfont\fontsize{13}{16}\selectfont\bfseries}{}{0pt}{}
\titleformat{\subsection}{\normalfont\bfseries}{}{0pt}{}
\titlespacing*{\section}{0pt}{1.4em}{0.6em}
\titlespacing*{\subsection}{0pt}{1em}{0.4em}

% --- Leyendas: "Cuadro 1." arriba y "Figura 1." abajo, en negrita y centradas -
\captionsetup{font=bf, labelsep=period, justification=centering}
\captionsetup[table]{position=top}

\begin{document}

% !TEX root = ../tecnologia-en-marcha.tex
% !TEX encoding = UTF-8 Unicode
% =============================================================================
% Título, autores, resumen y palabras clave.
% Texto: primer borrador del artículo (articulo_tecnologia_en_marcha.md), con los
% acentos restaurados.
% =============================================================================

% El título agrega lo que identifica a la investigación (PIA02): la ENIA y las
% personas adultas mayores. La traducción al inglés es una propuesta: la
% plantilla pide evitar traductores automáticos.
\newcommand{\tituloES}{Marco operativo de cumplimiento de la Estrategia Nacional de IA de Costa Rica para el monitoreo domiciliario de personas adultas mayores con sensores PIR}
\newcommand{\tituloEN}{Operational compliance framework for Costa Rica's National AI Strategy in home monitoring of older adults with PIR sensors}

% Una persona autora por bloque: nombre en negrita; afiliación y país; correo;
% ORCID. La plantilla deja lugar para tres.
\newcommand{\autor}[4]{%
  \textbf{#1}\\
  #2\\
  #3\\
  ORCID: #4\par\medskip}

\newcommand{\bloqueAutores}{%
  \autor{Miguel Ángel Méndez Arias}
        {Universidad CENFOTEC. Costa Rica}
        {\pendiente{correo}}
        {\pendiente{ORCID}}
  \autor{Laura María Segura Miranda}
        {Universidad CENFOTEC. Costa Rica}
        {\pendiente{correo}}
        {\pendiente{ORCID}}
  \pendiente{orden de las personas autoras y si hay una tercera (p. ej., quien dirigió la tesis)}\par}

\newcommand{\resumenES}{%
La Estrategia Nacional de Inteligencia Artificial de Costa Rica establece
principios de transparencia, equidad y responsabilidad, pero no especifica
mecanismos técnicos para verificar su cumplimiento en sistemas concretos. Este
artículo presenta y evalúa un marco operativo que traduce esos principios en
una batería de pruebas computacionales y artefactos auditables, aplicada a
sistemas de clasificación de actividades basados en sensores de infrarrojo
pasivo y otros sensores ambientales para el monitoreo domiciliario de personas
adultas mayores. La metodología consistió en traducir tres
principios de la estrategia nacional en nueve requerimientos operativos,
articularlos con el AI Act de la Unión Europea y el NIST AI Risk Management
Framework, e implementarlos como un procedimiento reproducible de seis pasos. El
marco integra explicabilidad local y global mediante SHAP, evaluación de equidad
por subgrupos con umbrales declarados previamente, registro estructurado de
inferencias y generación de artefactos auditables. La aplicación piloto sobre
dos configuraciones del conjunto CASAS produjo expedientes completos: Aruba, con
11\,175 inferencias, y una serie multi-hogar, con 7561 inferencias. En ambos
casos se cubrieron los nueve requerimientos y se obtuvo un veredicto de no
cumplimiento del sistema evaluado. Los resultados muestran que el marco permite
transformar principios normativos en evidencia técnica verificable y
reproducible.}

\newcommand{\palabrasclaveES}{%
Inteligencia artificial; sensores; aprendizaje automático; explicabilidad;
equidad; trazabilidad; auditoría
\revisar{la plantilla pide términos del tesauro de la UNESCO}}

\newcommand{\resumenEN}{%
Costa Rica's National Artificial Intelligence Strategy establishes principles
of transparency, fairness, and accountability, but does not specify technical
mechanisms to verify compliance with them in concrete systems. This article
presents and evaluates an operational framework that translates these
principles into a battery of computational tests and auditable artifacts,
applied to activity classification systems based on passive infrared and other
environmental sensors for home monitoring of older adults. The
methodology translated three principles of the national strategy into nine
operational requirements, aligned them with the European Union AI Act and the
NIST AI Risk Management Framework, and implemented them as a reproducible
six-step procedure. The framework integrates local and global explainability
through SHAP, subgroup fairness evaluation with thresholds declared in advance,
structured inference logging, and generation of auditable artifacts. The pilot
application on two CASAS dataset configurations produced complete evidence
files: Aruba, with 11,175 inferences, and a multi-home series, with 7,561
inferences. In both cases, all nine requirements were covered and the evaluated
system received a non-compliance verdict. The results show that the framework
can transform normative principles into verifiable and reproducible technical
evidence.}

\newcommand{\palabrasclaveEN}{%
Artificial intelligence; sensors; machine learning; explainability; fairness;
traceability; audit}


% --- Bloque de título ---------------------------------------------------------
\begin{center}
  {\LARGE\bfseries \tituloES\par}
  \vspace{0.6em}
  {\Large\bfseries \tituloEN\par}
  \vspace{1.4em}
\end{center}

{\setlength{\parindent}{0pt}\bloqueAutores}

% --- Resumen y abstract ----------------------------------------------------------
\section{Resumen}
\resumenES

\section{Palabras clave}
\palabrasclaveES

\section{Abstract}
\resumenEN

\section{Keywords}
\palabrasclaveEN

% --- Cuerpo ------------------------------------------------------------------
% !TEX root = ../tecnologia-en-marcha.tex
% !TEX encoding = UTF-8 Unicode
% Texto: borrador de Laura Segura, con acentos restaurados y citas enlazadas.

\section{Introducción}

Los sensores ambientales instalados en viviendas han sido utilizados para
reconocer actividades de la vida diaria sin recurrir a cámaras ni micrófonos
\cite{kim2022inhome,requena2024independence}. En particular, los sensores de infrarrojo pasivo, junto con sensores magnéticos de
puerta y sensores de temperatura, permiten registrar patrones de movimiento
compatibles con tareas como dormir, cocinar, comer o desplazarse dentro del
hogar. Esta clase de sistemas resulta relevante para el monitoreo no intrusivo
de personas adultas mayores, pues puede aportar información de apoyo a
cuidadores y servicios de acompañamiento \cite{park2025aging}. Sin embargo,
cuando las inferencias automatizadas se usan sobre poblaciones vulnerables, el
desempeño técnico del clasificador no es suficiente para justificar su uso.
% Las dos afirmaciones siguientes resumen la lectura que hace PIA02 (1.9 y
% 2.3.2) de esas revisiones.
Las revisiones recientes señalan, además, que la mayoría de los estudios del
campo se concentran en demostrar el desempeño técnico de los modelos y no en
evaluar su transparencia, su equidad o su gobernanza
\cite{park2025aging,habibovic2025aging}.

La Estrategia Nacional de Inteligencia Artificial de Costa Rica plantea
principios rectores para el desarrollo y uso responsable de sistemas de IA,
entre ellos transparencia y explicabilidad, equidad y no discriminación, y
responsabilidad \cite{micitt2024enia}. Estos principios son coherentes con
instrumentos internacionales como el Reglamento Europeo de Inteligencia
Artificial, que establece obligaciones de transparencia, gobernanza de datos y
conservación de registros para sistemas de alto riesgo \cite{ue2024aiact}, y con
el NIST AI Risk Management Framework, que organiza la gestión de riesgos en
funciones de mapear, medir, gobernar y gestionar \cite{nist2023airmf}. No
obstante, estos marcos suelen expresarse en un nivel normativo o metodológico
que no siempre indica cómo producir evidencia técnica verificable sobre un
sistema específico.

El problema abordado en este estudio es la brecha entre principios regulatorios
de alto nivel y mecanismos computacionales reproducibles para auditarlos en
sistemas de IA basados en sensores domiciliarios. En este dominio, la brecha es
particularmente relevante porque los datos describen rutinas domésticas
sensibles, los metadatos demográficos disponibles son limitados y los errores
pueden afectar de forma desigual a personas con patrones de movilidad o rutina
menos representados.

El objetivo de la investigación fue diseñar, implementar y verificar un marco
operativo que traduzca tres principios de la Estrategia Nacional de IA de Costa
Rica en requerimientos técnicos, pruebas computacionales y artefactos auditables
aplicables a un clasificador de actividades basado en sensores PIR. El artículo
resume la arquitectura del marco, la metodología de implementación y los
resultados obtenidos en dos aplicaciones piloto sobre datos públicos del
proyecto CASAS.

% Texto: borrador de Laura Segura, con acentos restaurados y citas enlazadas.

\section{Materiales y métodos}

\subsection{Tipo de investigación y enfoque}

La investigación fue de tipo aplicada y de diseño tecnológico. El objeto
evaluado no fue la optimización predictiva del modelo, sino la capacidad del
marco para generar evidencia verificable de cumplimiento. Por ello, el modelo de
aprendizaje automático se trató como sujeto de prueba: se entrenó y selló antes
de ejecutar el procedimiento de cumplimiento, y no se ajustó posteriormente con
base en los resultados de equidad o explicabilidad.

La metodología se organizó en cuatro componentes. El primero fue una matriz de
mapeo principio-requerimiento. El segundo correspondió a tres módulos de
cumplimiento: explicabilidad, equidad y trazabilidad. El tercero fue un registro
estructurado de inferencias. El cuarto fue un procedimiento reproducible de seis
pasos para generar el expediente de evidencia.

\subsection{Datos y población de estudio}

Se utilizaron datos públicos del proyecto CASAS, compuesto por eventos de
sensores ambientales en entornos domiciliarios \cite{cook2012aruba}. La primera
configuración corresponde al hogar Aruba, con una vivienda, sensores PIR,
sensores de puerta, sensores de temperatura y actividades anotadas. La segunda
corresponde a una serie multi-hogar, compuesta por nueve viviendas de personas
adultas mayores de una comunidad de retiro.
% El depósito de Zenodo de la serie hh pide citar Cook et al. (2013)
% [cook2013casas] y el propio conjunto [cook2025zenodo]; están en el .bib.
\pendiente{citar el depósito de la serie multi-hogar (Zenodo) y el artículo que pide citar}
En ambos casos, la población observada corresponde a personas adultas mayores
en viviendas independientes, aunque el proyecto no recolectó datos nuevos ni
tuvo contacto con participantes humanos.
% config.yaml describe a la residente de Aruba como "mujer adulta voluntaria":
% no dice que sea adulta mayor. Para la serie hh sí consta (comunidad de retiro).
\revisar{para Aruba, la documentación del conjunto dice ``mujer adulta'', no adulta mayor}

Los eventos se transformaron en ventanas disjuntas de 30 eventos consecutivos.
Cada ventana generó una fila tabular con conteos por zona del hogar, eventos de
puerta, variables temporales y, cuando estaban disponibles, lecturas de
temperatura. La partición de entrenamiento y prueba fue temporal por día, con
una proporción de prueba de 0,2. Esta decisión evitó mezclar ventanas contiguas
de una misma secuencia entre entrenamiento y prueba.

\subsection{Modelo de referencia}

El modelo de referencia fue un clasificador \texttt{RandomForestClassifier} con
semilla fija, ejecución de un solo hilo y versiones de bibliotecas declaradas.
Su función fue proporcionar inferencias para auditar el marco. La calidad
predictiva del modelo no se usó como criterio de éxito del estudio, pues el
propósito fue comprobar si el procedimiento producía evidencia completa y
reconstruible.
% Disponibles en el .bib: breiman2001rf (bosques aleatorios) y
% pedregosa2011sklearn (scikit-learn).

\subsection{Matriz de requerimientos}

La matriz tradujo los principios seleccionados en nueve requerimientos
operativos. El Cuadro~\ref{cuadro:requerimientos} resume la correspondencia
entre principio, requerimiento y evidencia esperada.

\begin{table}[htbp]
  \caption{Requerimientos operativos y evidencia del marco.}
  \label{cuadro:requerimientos}
  \small
  \begin{tabularx}{\linewidth}{@{}>{\raggedright\arraybackslash}p{0.27\linewidth}c>{\raggedright\arraybackslash}X@{}}
    \toprule
    Principio & Código & Evidencia técnica y documental \\
    \midrule
    Transparencia y explicabilidad & R3.1 & Explicación local por inferencia mediante atribución cuantificada de variables. \\
    Transparencia y explicabilidad & R3.2 & Importancia global agregada sobre el conjunto de evaluación. \\
    Transparencia y explicabilidad & R3.3 & Traducción de resultados a lenguaje comprensible y model card. \\
    \midrule
    Equidad y no discriminación & R4.1 & Desempeño desagregado por subgrupo. \\
    Equidad y no discriminación & R4.2 & Métricas de disparidad con umbrales declarados antes de medir. \\
    Equidad y no discriminación & R4.3 & Datasheet del conjunto de datos con procedencia, composición y limitaciones. \\
    \midrule
    Responsabilidad & R5.1 & Bitácora estructurada y persistente de inferencias. \\
    Responsabilidad & R5.2 & Model card con propósito, desempeño, limitaciones y condiciones de uso. \\
    Responsabilidad & R5.3 & Verificación de reconstrucción y atribución ex post de decisiones. \\
    \bottomrule
  \end{tabularx}
  \fuente{elaboración propia.}
\end{table}

\subsection{Módulo de explicabilidad}

El módulo de explicabilidad utilizó SHAP con TreeExplainer para calcular
atribuciones locales de las variables de entrada
\cite{lundberg2017shap,lundberg2020treeshap}. Para cada inferencia del conjunto
de prueba se registró la clase predicha, el valor base, las contribuciones de
las variables y una referencia al artefacto de explicación. La importancia
global se obtuvo como la media de las contribuciones absolutas locales, evitando
un segundo cálculo independiente y manteniendo coherencia entre explicaciones
locales y globales. Las variables se nombraron por zonas del hogar, por ejemplo
cocina o dormitorio, para favorecer la comprensión por parte de destinatarios no
técnicos.

\subsection{Módulo de equidad}

El módulo de equidad evaluó desempeño desagregado y disparidad entre subgrupos.
En Aruba, los subgrupos fueron contextuales, como franja horaria y tipo de día.
En la configuración multi-hogar, se incorporó el hogar como subgrupo
poblacional, pues cada vivienda representa una persona o unidad doméstica
distinta. Las métricas con efecto en el veredicto fueron diferencia de tasa de
verdaderos positivos y diferencia de tasa de falsos positivos, ambas con umbral
de 0,10 declarado antes de la ejecución. También se reportaron como descriptivas
la diferencia de paridad demográfica y el cociente de tasa de selección.
% Disponible en el .bib: hardt2016equality, origen de la igualdad de
% oportunidades (diferencia de TVP).

\subsection{Módulo de trazabilidad}

La trazabilidad se implementó mediante una bitácora JSONL \emph{append-only}.
Cada línea registró identificador de evento, marca temporal UTC, hash de la
entrada procesada, salida del modelo, confianza, versión del modelo, versión del
marco, referencia a la explicación local, responsable declarado y versión del
esquema. El uso de hashes evitó persistir datos crudos sensibles en la
bitácora, sin impedir la reconstrucción técnica de la decisión dentro del
expediente.

\subsection{Procedimiento de aplicación}

El procedimiento siguió seis pasos alineados con las funciones del NIST AI RMF:
caracterización del sistema, documentación del conjunto de datos, declaración de
criterios de evaluación, ejecución de pruebas técnicas, generación de artefactos
auditables y verificación de auditabilidad. El paso de declaración de criterios
selló el hash SHA-256 del archivo de configuración. En la verificación final, el
marco comprobó que ese hash no hubiera cambiado, que la bitácora fuera válida,
que existiera cobertura uno a uno entre inferencias esperadas y registradas, y
que los nueve requerimientos contaran con artefactos presentes.
% Disponibles en el .bib: gebru2021datasheets y mitchell2019modelcards, que
% originan el datasheet y el model card del procedimiento.

% !TEX root = ../tecnologia-en-marcha.tex
% !TEX encoding = UTF-8 Unicode
% Texto: primer borrador del artículo, con acentos restaurados y citas enlazadas.

\section{Resultados}

El marco fue ejecutado sobre dos configuraciones piloto. El
Cuadro~\ref{cuadro:resultados} sintetiza los resultados generales del
expediente.

\begin{table}[htbp]
  \caption{Resultados generales de las aplicaciones piloto.}
  \label{cuadro:resultados}
  \footnotesize
  % Las columnas numéricas son p{} alineadas a la derecha, no r: así el
  % encabezado ("Inferencias registradas") puede partirse en dos líneas.
  % TeX no corta con guion la primera palabra de un párrafo, así que cada
  % columna tiene que caber su palabra más larga ("Requerimientos").
  \begin{tabularx}{\linewidth}{@{}%
      >{\raggedright\arraybackslash}X%
      >{\raggedright\arraybackslash}p{0.23\linewidth}%
      >{\raggedleft\arraybackslash}p{0.10\linewidth}%
      >{\raggedleft\arraybackslash}p{0.12\linewidth}%
      >{\raggedleft\arraybackslash}p{0.15\linewidth}%
      >{\raggedright\arraybackslash}p{0.12\linewidth}@{}}
    \toprule
    Piloto & Configuración & Viviendas & Inferencias registradas & Requerimientos cubiertos & Veredicto del sistema evaluado \\
    \midrule
    CASAS Aruba    & \texttt{config.yaml}         & 1 & 11\,175 & 9 de 9 & No cumple \\
    CASAS serie hh & \texttt{config.hogares.yaml} & 9 & 7561    & 9 de 9 & No cumple \\
    \bottomrule
  \end{tabularx}
  \fuente{elaboración propia con base en la ejecución documentada del repositorio.}
\end{table}

Los dos expedientes generaron la totalidad de los artefactos previstos: ficha de
caracterización, datasheet, protocolo de evaluación, explicaciones locales,
importancia global, resultados de desempeño desagregado, métricas de equidad,
bitácora de inferencias, reporte de explicabilidad, reporte de equidad, model
card, reporte de cumplimiento, bitácora de ejecución y manifiesto de hashes. En
ambos pilotos, los nueve requerimientos de la matriz quedaron respaldados por
evidencia presente y verificable.
% El manifiesto sella 14 artefactos: los de esta lista menos el propio
% manifiesto, más la figura de importancia global.

El resultado de no cumplimiento no indica un fallo del marco, sino un hallazgo
sobre el sistema evaluado. El procedimiento completo se ejecutó y produjo
evidencia suficiente para sostener el veredicto. Esto es consistente con el
criterio de éxito definido para el piloto: el marco debía demostrar que podía
generar un expediente de evidencia, incluso cuando la evidencia revelara
incumplimientos del sistema examinado.

En el módulo de explicabilidad, cada inferencia registrada en la bitácora quedó
vinculada con una explicación local. Esta relación satisface simultáneamente
R3.1 y R5.3, porque permite partir de una decisión registrada y llegar al
artefacto que explica las variables que contribuyeron a la predicción. La
agregación de esas explicaciones permitió construir una importancia global de
variables, asociada con R3.2. El uso de nombres de zonas del hogar permitió
generar enunciados interpretables para cuidadores y otros destinatarios no
técnicos, en concordancia con R3.3.

% Fuente: artefactos/reporte_explicabilidad.md y
% artefactos/hogares/reporte_explicabilidad.md (media de |SHAP|).
En Aruba, las variables con mayor importancia global, medida como la media del
valor absoluto de SHAP, fueron el movimiento en el sillón (0,0887), en la cocina
(0,0780) y en el dormitorio (0,0447), seguidas de la duración de la ventana y la
hora del día. En la configuración multi-hogar encabezaron la cocina (0,0452), el
baño (0,0440) y el dormitorio (0,0406). En ambos casos, las variables
dominantes son zonas del hogar, lo que permite contrastar el comportamiento del
modelo con el conocimiento del dominio sobre dónde ocurre cada actividad.

En el módulo de equidad, los resultados fueron desagregados por los subgrupos
declarados en la configuración. La existencia de umbrales sellados antes de
ejecutar las pruebas permitió diferenciar la medición de disparidades de una
interpretación posterior ajustada a los resultados. Cuando una combinación de
actividad, subgrupo y métrica no tuvo soporte suficiente, el marco la marcó como
no evaluable en lugar de contarla como aprobada. Esta decisión evitó convertir
ausencia de evidencia en evidencia de cumplimiento.

% Todas las cifras de esta parte salen de los expedientes: model_card.md,
% reporte_equidad.md y equidad/desempeno_desagregado.csv de cada piloto.
Como referencia para interpretar las disparidades, el modelo alcanzó una
exactitud de 0,704 y un F1 macro de 0,432 en Aruba, y de 0,326 y 0,276 en la
configuración multi-hogar, con una exactitud por vivienda entre 0,225 y 0,433.
Estas cifras no forman parte del criterio de éxito del estudio, pero condicionan
la lectura de los resultados de equidad: un clasificador de desempeño modesto
deja más margen para que aparezcan diferencias entre subgrupos.

En Aruba se evaluaron 32 de las 44 combinaciones de subgrupo, actividad y
métrica con efecto en el veredicto; las 12 restantes no alcanzaron el soporte
mínimo de 30 casos por tasa. Cuatro combinaciones superaron el umbral
(Cuadro~\ref{cuadro:aruba}). Las cuatro corresponden a la franja horaria y a la
diferencia de TVP: el modelo reconoció peor ciertas actividades según la hora
del día, mientras que ninguna combinación de tasa de falsos positivos ni de tipo
de día superó el umbral. Parte de esta disparidad puede deberse al diseño, pues
la franja se deriva de la hora, que es también una variable de entrada del
modelo. La mayor diferencia, en Sleeping, compara solo la madrugada y la noche,
porque la mañana y la tarde quedaron por debajo del soporte mínimo.

\begin{table}[htbp]
  \caption{Combinaciones que superan el umbral de 0,10 en Aruba.}
  \label{cuadro:aruba}
  \small
  \begin{tabularx}{\linewidth}{@{}>{\raggedright\arraybackslash}X*{5}{r}@{}}
    \toprule
    & & \multicolumn{4}{c}{TVP por franja horaria} \\
    \cmidrule(l){3-6}
    Actividad & Diferencia de TVP & Madrugada & Mañana & Tarde & Noche \\
    \midrule
    Sleeping         & 0,351 & 0,970 & 0,944* & 0,500* & 0,619 \\
    Otro             & 0,170 & 0,522 & 0,692  & 0,530  & 0,622 \\
    Meal\_Preparation & 0,169 & 0,880 & 0,870  & 0,798  & 0,712 \\
    Relax            & 0,128 & 0,835 & 0,901  & 0,773  & 0,820 \\
    \bottomrule
  \end{tabularx}
  \par\smallskip{\footnotesize * Menos de 30 casos: fuera de la comparación.\par}
  \fuente{reporte de equidad del expediente de Aruba.}
\end{table}

Entre las combinaciones no evaluables está Bed\_to\_Toilet, los traslados
nocturnos al baño: la prueba contiene 8 casos en la madrugada y ninguno en las
demás franjas. El marco no puede afirmar que el sistema reconozca por igual
esta actividad; sí verificó que sus falsas alarmas se mantuvieran dentro del
umbral en todas las franjas.

En la configuración multi-hogar se evaluaron 108 de 150 combinaciones y 29
superaron el umbral: 14 entre viviendas, 12 entre franjas horarias y 3 entre
tipos de día. A diferencia de Aruba, aquí los subgrupos corresponden a personas
distintas y no a condiciones de contexto, como plantea el requerimiento R4.1.
Las mayores diferencias entre viviendas se dieron en la TVP de Work (0,872),
Eat (0,810), Cook (0,503), Wash\_Dishes (0,447) y Sleep (0,425).

El caso de Cook muestra que la disparidad no se explica solo por la cantidad de
datos. La vivienda hh107, la de más casos de esa actividad en la prueba (243),
obtuvo la TVP más baja (0,292), mientras que hh103, con 122, obtuvo la más alta
(0,795). En Sleep, en cambio, hh107 superó a hh104 (0,867 frente a 0,442), las
dos únicas viviendas con soporte suficiente para esa actividad. El desempeño
desigual depende, así, de la combinación entre vivienda y actividad.
% D52 atribuye el caso de hh107 a que es la única vivienda con dos residentes.
% El dato no consta en el expediente ni en el crudo del repositorio.
\revisar{D52 atribuye la diferencia de hh107 a que tiene dos residentes; hace falta la fuente de ese dato para incluirlo}

En el módulo de trazabilidad, la bitácora estructurada permitió verificar
unicidad de identificadores, marcas temporales UTC, rango válido de confianza,
versión de modelo, versión de esquema, responsable declarado y existencia de la
explicación referenciada. La verificación final agregó la cobertura entre
conjunto de prueba y registros producidos, de modo que una bitácora válida pero
incompleta no pudiera sostener el cumplimiento.

Adicionalmente, el repositorio documentó 272 pruebas automatizadas en verde.
Estas pruebas cubrieron configuración, preparación de datos, modelado,
explicabilidad, equidad, trazabilidad, procedimiento y documentación. Aunque las
pruebas no sustituyen una auditoría externa, aportan evidencia de
reproducibilidad del procedimiento y reducen el riesgo de divergencia entre el
diseño normativo y la implementación computacional.

% Texto: borrador de Laura Segura, con acentos restaurados y citas enlazadas.

\section{Discusión}

Los resultados muestran que es posible operacionalizar principios regulatorios
mediante requerimientos discretos, pruebas computacionales y artefactos
auditables. La principal contribución del marco es que no se limita a declarar
alineación con la Estrategia Nacional de IA, el AI Act o el NIST AI RMF, sino
que produce evidencias verificables: hashes de configuración, registros
estructurados, reportes, model card, datasheet y matriz de cobertura.

La explicabilidad basada en SHAP resultó adecuada para un modelo de árboles y
datos tabulares. Su ventaja principal fue conectar cada inferencia con
atribuciones cuantificadas y, posteriormente, agregar esas atribuciones para
describir el comportamiento global. No obstante, esta selección también impone
una limitación: el desempeño computacional y la validez de las explicaciones
dependen de la familia de modelos y del modo de perturbación seleccionado.

La evaluación de equidad evidenció una restricción importante del dominio. Los
conjuntos públicos de sensores domiciliarios suelen incluir pocos metadatos
demográficos, por lo que no siempre permiten comparar categorías protegidas
como género, etnia, capacidad económica o nivel formativo. El piloto
multi-hogar mitigó parcialmente esta restricción al comparar desempeño entre
viviendas, pero no la elimina por completo. En consecuencia, el marco debe
entenderse como una herramienta para producir evidencia disponible y declarar
sus límites, no como garantía absoluta de ausencia de discriminación.
% Límites que los expedientes declaran y el borrador todavía no menciona:
%   - la franja horaria sale de la hora, que es también variable de entrada:
%     parte de la disparidad entre franjas es por diseño (config.yaml);
%   - la exactitud global del piloto multi-hogar es 0,22--0,43 por vivienda,
%     contra 0,70 en Aruba; un modelo más débil tiene más lugar para mostrar
%     disparidad (D52);
%   - no se midió robustez ni la varianza entre semillas (D56).
\pendiente{decidir si se incluyen los límites que declara el expediente: la hora como variable de entrada, la exactitud baja del multi-hogar, robustez y varianza entre semillas}

La trazabilidad fue el componente transversal que permitió unir resultados
técnicos y responsabilidad. Sin una bitácora persistente, las explicaciones y
los reportes quedarían desconectados de decisiones concretas. Con el esquema
propuesto, cada decisión puede reconstruirse a partir de su identificador,
entrada hasheada, salida, confianza, versión del modelo, explicación y
responsable declarado.

Una diferencia relevante frente a procesos generales de aseguramiento de calidad
en aprendizaje automático, como CRISP-ML(Q), es que el marco no busca optimizar
el modelo ni seleccionar la mejor arquitectura \cite{studer2021crispmlq}. Su
propósito es auditar un sistema dado. Por ello, un veredicto de no cumplimiento
es un resultado válido: revela que el sistema evaluado no satisface los
criterios declarados, pero también demuestra que el procedimiento puede
producir evidencia para sostener esa conclusión.

% !TEX root = ../tecnologia-en-marcha.tex
% !TEX encoding = UTF-8 Unicode
% Texto: primer borrador del artículo, con acentos restaurados.

\section{Conclusiones y recomendaciones}

El marco operativo desarrollado permitió traducir tres principios de la
Estrategia Nacional de IA de Costa Rica en nueve requerimientos verificables
para sistemas de IA basados en sensores PIR. La implementación demostró que esos
requerimientos pueden evaluarse mediante un procedimiento reproducible de seis
pasos, articulado con el AI Act de la Unión Europea y el NIST AI RMF.

La aplicación sobre dos pilotos del conjunto CASAS generó expedientes
completos, con 18\,736 inferencias registradas en total y cobertura de los
nueve requerimientos en ambos casos. Los dos sistemas evaluados obtuvieron
veredicto de no cumplimiento, lo cual confirma la utilidad del marco para
producir evidencia incluso cuando el resultado es desfavorable para el sistema
auditado.

El estudio aporta una forma concreta de pasar de principios normativos a
evidencia técnica. En particular, integra explicabilidad local y global,
medición de equidad con umbrales declarados previamente, trazabilidad por
inferencia y documentos auditables. Esta combinación permite que un tercero
revise no solo los resultados, sino también el protocolo y los artefactos que
los sostienen.

Como trabajo futuro, se recomienda validar los umbrales de equidad con un panel
multidisciplinario, incorporar pruebas de robustez ante entradas ruidosas o
degradadas, evaluar reproducibilidad de resultado sobre múltiples semillas y
extender la verificación a escenarios de monitoreo continuo. También se
recomienda aplicar el marco a datasets con metadatos demográficos más
completos, siempre que su uso sea éticamente justificable y compatible con la
protección de datos personales.

% Texto: borrador de Laura Segura, con acentos restaurados. La segunda frase de
% agradecimientos era una nota para las personas autoras: pasa a \pendiente.

\section{Agradecimientos}

Se agradece al proyecto CASAS de Washington State University por la
disponibilidad pública de conjuntos de datos de sensores ambientales para
investigación en reconocimiento de actividades.
\pendiente{apoyos institucionales, financiamiento o contribuciones específicas, si corresponde}

\section{Declaración sobre el uso de inteligencia artificial}

Los autores declaramos que hemos utilizado GitHub Copilot para asistir en la
estructuración inicial y redacción preliminar de este borrador de artículo. La
herramienta ayudó a organizar el texto conforme a las instrucciones editoriales
y a mejorar la claridad expositiva. Los contenidos, cifras, referencias y
conclusiones deben ser revisados, corregidos y validados por los autores antes
de cualquier envío, quienes conservan la responsabilidad intelectual y ética
sobre la versión final.
% La conversión a LaTeX y la verificación de referencias se hicieron con
% Claude (Anthropic). Si la revista lo exige, la declaración debe decirlo.
\revisar{la conversión a LaTeX y la verificación de referencias se hicieron con Claude; actualizar la declaración}


\printbibliography[heading=bibnumbered,title={Referencias}]

\end{document}
